% Options for packages loaded elsewhere
% Options for packages loaded elsewhere
\PassOptionsToPackage{unicode}{hyperref}
\PassOptionsToPackage{hyphens}{url}
%
\documentclass[
  ignorenonframetext,
]{beamer}
\newif\ifbibliography
\usepackage{pgfpages}
\setbeamertemplate{caption}[numbered]
\setbeamertemplate{caption label separator}{: }
\setbeamercolor{caption name}{fg=normal text.fg}
\beamertemplatenavigationsymbolsempty
% remove section numbering
\setbeamertemplate{part page}{
  \centering
  \begin{beamercolorbox}[sep=16pt,center]{part title}
    \usebeamerfont{part title}\insertpart\par
  \end{beamercolorbox}
}
\setbeamertemplate{section page}{
  \centering
  \begin{beamercolorbox}[sep=12pt,center]{section title}
    \usebeamerfont{section title}\insertsection\par
  \end{beamercolorbox}
}
\setbeamertemplate{subsection page}{
  \centering
  \begin{beamercolorbox}[sep=8pt,center]{subsection title}
    \usebeamerfont{subsection title}\insertsubsection\par
  \end{beamercolorbox}
}
% Prevent slide breaks in the middle of a paragraph
\widowpenalties 1 10000
\raggedbottom
\AtBeginPart{
  \frame{\partpage}
}
\AtBeginSection{
  \ifbibliography
  \else
    \frame{\sectionpage}
  \fi
}
\AtBeginSubsection{
  \frame{\subsectionpage}
}
\usepackage{iftex}
\ifPDFTeX
  \usepackage[T1]{fontenc}
  \usepackage[utf8]{inputenc}
  \usepackage{textcomp} % provide euro and other symbols
\else % if luatex or xetex
  \usepackage{unicode-math} % this also loads fontspec
  \defaultfontfeatures{Scale=MatchLowercase}
  \defaultfontfeatures[\rmfamily]{Ligatures=TeX,Scale=1}
\fi
\usepackage{lmodern}

\usetheme[]{metropolis}
\ifPDFTeX\else
  % xetex/luatex font selection
\fi
% Use upquote if available, for straight quotes in verbatim environments
\IfFileExists{upquote.sty}{\usepackage{upquote}}{}
\IfFileExists{microtype.sty}{% use microtype if available
  \usepackage[]{microtype}
  \UseMicrotypeSet[protrusion]{basicmath} % disable protrusion for tt fonts
}{}
\makeatletter
\@ifundefined{KOMAClassName}{% if non-KOMA class
  \IfFileExists{parskip.sty}{%
    \usepackage{parskip}
  }{% else
    \setlength{\parindent}{0pt}
    \setlength{\parskip}{6pt plus 2pt minus 1pt}}
}{% if KOMA class
  \KOMAoptions{parskip=half}}
\makeatother


\usepackage{longtable,booktabs,array}
\usepackage{calc} % for calculating minipage widths
\usepackage{caption}
% Make caption package work with longtable
\makeatletter
\def\fnum@table{\tablename~\thetable}
\makeatother
\usepackage{graphicx}
\makeatletter
\newsavebox\pandoc@box
\newcommand*\pandocbounded[1]{% scales image to fit in text height/width
  \sbox\pandoc@box{#1}%
  \Gscale@div\@tempa{\textheight}{\dimexpr\ht\pandoc@box+\dp\pandoc@box\relax}%
  \Gscale@div\@tempb{\linewidth}{\wd\pandoc@box}%
  \ifdim\@tempb\p@<\@tempa\p@\let\@tempa\@tempb\fi% select the smaller of both
  \ifdim\@tempa\p@<\p@\scalebox{\@tempa}{\usebox\pandoc@box}%
  \else\usebox{\pandoc@box}%
  \fi%
}
% Set default figure placement to htbp
\def\fps@figure{htbp}
\makeatother

\ifLuaTeX
  \usepackage{luacolor}
  \usepackage[soul]{lua-ul}
\else
  \usepackage{soul}
  \makeatletter
  \let\HL\hl
  \renewcommand\hl{% fix for beamer highlighting
    \let\set@color\beamerorig@set@color
    \let\reset@color\beamerorig@reset@color
    \HL}
  \makeatother
\fi




\setlength{\emergencystretch}{3em} % prevent overfull lines

\providecommand{\tightlist}{%
  \setlength{\itemsep}{0pt}\setlength{\parskip}{0pt}}



 


\setbeamerfont{title}{size=\large} \usepackage{hyperref} \setbeamertemplate{footline}[frame number]
\makeatletter
\@ifpackageloaded{caption}{}{\usepackage{caption}}
\AtBeginDocument{%
\ifdefined\contentsname
  \renewcommand*\contentsname{Table of contents}
\else
  \newcommand\contentsname{Table of contents}
\fi
\ifdefined\listfigurename
  \renewcommand*\listfigurename{List of Figures}
\else
  \newcommand\listfigurename{List of Figures}
\fi
\ifdefined\listtablename
  \renewcommand*\listtablename{List of Tables}
\else
  \newcommand\listtablename{List of Tables}
\fi
\ifdefined\figurename
  \renewcommand*\figurename{Figure}
\else
  \newcommand\figurename{Figure}
\fi
\ifdefined\tablename
  \renewcommand*\tablename{Table}
\else
  \newcommand\tablename{Table}
\fi
}
\@ifpackageloaded{float}{}{\usepackage{float}}
\floatstyle{ruled}
\@ifundefined{c@chapter}{\newfloat{codelisting}{h}{lop}}{\newfloat{codelisting}{h}{lop}[chapter]}
\floatname{codelisting}{Listing}
\newcommand*\listoflistings{\listof{codelisting}{List of Listings}}
\makeatother
\makeatletter
\makeatother
\makeatletter
\@ifpackageloaded{caption}{}{\usepackage{caption}}
\@ifpackageloaded{subcaption}{}{\usepackage{subcaption}}
\makeatother

\usepackage{bookmark}
\IfFileExists{xurl.sty}{\usepackage{xurl}}{} % add URL line breaks if available
\urlstyle{same}
\hypersetup{
  pdfauthor={Giorgio Arcara},
  hidelinks,
  pdfcreator={LaTeX via pandoc}}


\author{Giorgio Arcara}
\date{}

\begin{document}


\begin{frame}
% TITLE SLIDES
\title{Metodologia della Ricerca Clinica\\ \vspace{1em} \emph{1. Introduzione}}
\author{Giorgio Arcara,\\ Università di Padova \\ IRCCS San Camillo, Venezia}

\titlegraphic{

\vspace*{7cm}
\includegraphics[scale=0.15]{Figures/LogoSanCamilloIRCCS_Unipd_alpha.png}
\begin{center}
\vspace{-2.2em}
\includegraphics[scale=0.15]{Figures/CC_license_3_0.png}
\end{center}
}
\date{\today}
\maketitle
\end{frame}

\begin{frame}{Lezione di oggi}
\phantomsection\label{lezione-di-oggi}
\begin{itemize}
\tightlist
\item
  Presentazione del docente e del corso.
\item
  Aspetti organizzativi del corso.
\item
  Obiettivi formativi e ``spirito'' del corso.
\end{itemize}
\end{frame}

\begin{frame}{CV Giorgio Arcara}
\phantomsection\label{cv-giorgio-arcara}
\footnotesize

\begin{itemize}
\item
  Laurea in Psicologia presso Unipd (2005)
\item
  Master in Neuropsicologia dei disturbi cognitivi acquisiti (2006)
\item
  Vari post-doc presso dipartimenti (psicologia, ingegneria,
  neuroscienze)
\item
  Dal 2016 Ricercatore presso IRCCS San Camillo di Venezia.
\item
  Dal 2025 Professore Associato in Psicometria - Università di Padova.
\end{itemize}

\pause

\hfill
\includegraphics[width=0.6\linewidth,height=\textheight,keepaspectratio]{Figures/GEMS.png}
\end{frame}

\begin{frame}{Contatti e riferimenti}
\phantomsection\label{contatti-e-riferimenti}
\centering

\href{https://giorgioarcara.github.io}{\ul{https://giorgioarcara.github.io}}

\href{https://www.dpg.unipd.it/category/ruoli/personale-docente?key=4D0707A2DD6FCCD50B90E4C9F316A625}{\ul{pagina
personale unipd}}

\href{mailto:giorgio.arcara@unipd.it}{\nolinkurl{giorgio.arcara@unipd.it}}
\end{frame}

\begin{frame}{Aspetti organizzativi - Materiali}
\phantomsection\label{aspetti-organizzativi---materiali}
\textbf{Condividerò tutto il materiale utilizzato e mostrato}

I materiali principali sono le slides e materiali aggiuntivi che vi saranno forniti durante il corso. I materiali li troverete anche nella pagina moodle.
\end{frame}

\begin{frame}{Aspetti organizzativi - Materiali}
\phantomsection\label{aspetti-organizzativi---materiali-1}
\begin{columns}
\column{0.5\textwidth}
\small
\emph{Mondini, S., Cappelletti, M., \& Arcara, G. (2022). Methodology in Neuropsychological Assessment: An Interpretative Approach to Guide Clinical Practice. Taylor \& Francis.}\\
\vspace{1em}
Testo di approfondimento, \textbf{non indispensabile}.

\column{0.5\textwidth}

\begin{figure}
\includegraphics[scale=0.5]{Figures/Methodology_book.png}
\end{figure}

\end{columns}
\end{frame}

\begin{frame}{Aspetti organizzativi - Materiali aggiuntivi}
\phantomsection\label{aspetti-organizzativi---materiali-aggiuntivi}
\begin{itemize}
\item
  Libro gratuito su statistica base ed R
  \href{https://learningstatisticswithr.com/}{\underline{https://learningstatisticswithr.com/}}
\item
  Libro gratuito su psicometria (più avanzato)
  \href{https://personality-project.org/r/book/}{\underline{https://personality-project.org/r/book/}}
\item
  Libro in corso di scrittura
  \href{https://giorgioarcara.github.io/oip-book/}{\underline{https://giorgioarcara.github.io/oip-book/}}
\end{itemize}
\end{frame}

\begin{frame}{Qualche domanda}
\phantomsection\label{qualche-domanda}
\emph{Chi tra di voi pensa che principalmente svolgerà attività
clinica?} \pause

\emph{Chi tra di voi pensa che principalmente si occuperà di ricerca?}
\pause

\emph{Chi tra di voi pensa che nella pratica clinica utilizzerà test?}
\end{frame}

\begin{frame}{L'approccio}
\phantomsection\label{lapproccio}
\centering

\textbf{Learning-by-doing}

\textbf{Active-learning}
\end{frame}

\begin{frame}{Clinica e ricerca}
\phantomsection\label{clinica-e-ricerca}
\begin{center}
\includegraphics[height=7cm]{Figures/Separazione.png}
\end{center}
\end{frame}

\begin{frame}{Obiettivi del corso}
\phantomsection\label{obiettivi-del-corso}
\begin{itemize}
\tightlist
\item
  Spiegare come utilizzare i test in clinica \pause
\item
  Spiegare i limiti dell'utilizzo dei test \pause
\item
  Spiegare fondamenti di ricerca clinica (per chi fa clinica)
\end{itemize}
\end{frame}

\begin{frame}{Un passo indietro per riflettere}
\phantomsection\label{un-passo-indietro-per-riflettere}
\begin{center}
\includegraphics[height=8cm]{Figures/Stepping_back.png}
\end{center}
\end{frame}

\begin{frame}{Utilizzo critico di LLM (es. ChatGPT) per il corso 1/2}
\phantomsection\label{utilizzo-critico-di-llm-es.-chatgpt-per-il-corso-12}
I Large Language Models (LLM) come ChatGPT, comunemente chiamati sistemi
di intelligenza artificiale sono una grande opportunità di innovazione
tecnologica che merita di essere discussa. Di seguito alcune
raccomandazioni per un utilizzo appropriato per il corso.
\end{frame}

\begin{frame}{Utilizzo critico di LLM (es. ChatGPT) per il corso 2/2}
\phantomsection\label{utilizzo-critico-di-llm-es.-chatgpt-per-il-corso-22}
\begin{itemize}
\item
  Non cercate risposte, ma cercate \emph{una guida} per trovare
  informazioni. Gli LLM funzionano come motori di ricerca avanzati
  soprattutto se la domanda è vaga.
\item
  Utilizzatele per capire meglio alcuni concetti. Se qualcosa non è
  chiaro, potete esplicitamente chiedere di spiegarlo in altri termini o
  con altri esempi.
\item
  \textbf{Non vi fidate ciecamente di ciò che dicono} ma cercate
  conferma da eventuali fonti. Gli LLM spesso mostrano dei comportamenti
  detti \emph{hallucinations} danno risposte verosimili o credibili ma
  che si rivelano sbagliate o completamente inventate.
\end{itemize}
\end{frame}




\end{document}
